{\renewcommand{\thefootnote}{(97)}\nde{One has added what follows.}}%
One can now complete the proof of~4.8\,(iii). The equality~\((\ast\ast\ast)\)
is checked locally on~\(T\), so one is reduced to the affine situation
described above.
Denote by~\(d_{A/\Lambda}\) the differential \(A\to\Omega^1_{A/\Lambda}\). Then
\(a\) corresponds, above~\(U\), to an element~\(a_U\) of
\[
\Hom_{C_0}(\Omega^1_{A_0/\Lambda_0}\otimes_{A_0}C_0,JC)
\simeq
\Hom_{B_0}(\Omega^1_{A/\Lambda}\otimes_A B_0,JC)
\simeq
\Hom_A(\Omega^1_{A/\Lambda},JC).
\]
Then, on the one hand, \(v_{f_0}(a)\) corresponds above~\(U\) to the
element \(a_U\circ D_0\), where \(D_0\) is the morphism of
\(B\)-modules%
{\renewcommand{\thefootnote}{(98)}\nde{\Cf N.D.E.~(93).}}
\[
I_Y/I_Y^2\longrightarrow\Omega^1_{A/\Lambda}\otimes_A B_0,
\qquad
x+I_Y^2\mapsto d_{A/\Lambda}(x)\otimes 1.
\]
On the other hand (\cf the proof of~0.1.8 and~0.1.9), the morphism of
\(\Lambda\)-algebras \(\phi'\colon A\to C\) corresponding to \(a\cdot f\)
differs from~\(\phi\) by the \(\Lambda\)-derivation \(A\to JC\) associated
with~\(a_U\), \ie one has:
\[
\phi'=\phi+a_U\circ d_{A/\Lambda}=\phi+a_U\circ(d_{A/\Lambda}\otimes 1).
\]
Consequently, writing \(c'=c(X,Y,a\cdot f)\), one has for every
\(x\in I_Y\), writing \(x\) for its image in \(I_Y/I_Y^2\):
\[
(c'-c)(x\otimes 1)
=a_U(d_{A/\Lambda}(x)\otimes 1)
=(a_U\circ D_0)(x)
=v_{f_0}(a)(x).
\]
This shows that \(c'-c=v_{f_0}(a)\).

\numpara{4.10}\label{III.4.10}
One has deleted remark~4.10 of the original, rendered obsolete by the
addition of remark~4.8.0.

\numpara{4.11}\label{III.4.11}
We now propose to study the following situation. Let \(S\),
\(S_{\mathcal{J}}\) and \(S_0\) be as in~4.8; one has three \(S\)-schemes
\(X\), \(X'\), \(T\), a subscheme \(Y\) of \(X\) (\resp \(Y'\) of \(X'\)), and
morphisms \(f\colon T\to X'\) and \(g\colon X'\to X\).
\[
\xymatrix@C=3.2em@R=1.6em{
& Y' \ar@{^{(}->}[d]^{i'} & Y \ar@{^{(}->}[d]^{i} \\
T \ar[r]^{f} & X' \ar[r]^{g} & X\,.
}
\]

One supposes that by reduction modulo~\(\mathcal{J}\), this diagram
completes to a commutative diagram
\[
\xymatrix@C=3.4em@R=1.7em{
& Y'_{\mathcal{J}} \ar@{-->}[rr] \ar@{^{(}->}[d]_{i'_{\mathcal{J}}}
& & Y_{\mathcal{J}} \ar@{^{(}->}[d]^{i_{\mathcal{J}}} \\
T_{\mathcal{J}} \ar@{-->}[ur] \ar[r]_{f_{\mathcal{J}}}
& X'_{\mathcal{J}} \ar[rr]_{g_{\mathcal{J}}}
& & X_{\mathcal{J}}\,.
}
\]
One therefore has, by~4.8, obstructions:
\begin{align*}
c(X,Y,g\circ i')
&\in
\Hom_{\mathcal{O}_{Y'_0}}\bigl(
i'^*_0 g^*_0(\mathcal{N}_{Y/X}\otimes_{\mathcal{O}_Y}\mathcal{O}_{Y_0}),
\mathcal{J}\mathcal{O}_{Y'}\bigr),\\
c(X',Y',f)
&\in
\Hom_{\mathcal{O}_{T_0}}\bigl(
f^*_0(\mathcal{N}_{Y'/X'}\otimes_{\mathcal{O}_{Y'}}\mathcal{O}_{Y'_0}),
\mathcal{J}\mathcal{O}_T\bigr),\\
c(X,Y,g\circ f)
&\in
\Hom_{\mathcal{O}_{T_0}}\bigl(
f^*_0 g^*_0(\mathcal{N}_{Y/X}\otimes_{\mathcal{O}_Y}\mathcal{O}_{Y_0}),
\mathcal{J}\mathcal{O}_T\bigr),
\end{align*}
whose relations one seeks to calculate.%
{\renewcommand{\thefootnote}{(99)}\nde{From~4.17 onwards, one will apply
this to the case where \(X\) is an \(S\)-group, \(g\colon X\times_S X\to X\)
the multiplication, \(Y\) a subscheme of \(X\) such that \(Y_{\mathcal{J}}\)
is a subgroup of \(X_{\mathcal{J}}\), \(Y'=Y\times_S Y\), and to the two
morphisms \(Y^3\to X^2\) which send \((y_1,y_2,y_3)\) to
\((y_1 y_2,y_3)\), \resp \((y_1,y_2 y_3)\). In this case, the comparison of
the obstructions above will show that the obstruction to \(Y\) being a
subgroup of \(X\) lies in a certain cohomology group (of Hochschild)
\(\mathrm{H}^2(Y_0,N_0)\).}}

\begin{lemma}{4.12}\label{III.4.12}
Suppose \(Y'\) flat over \(S\), so that
\(\mathcal{J}\otimes_{\mathcal{O}_{S_0}}\mathcal{O}_{Y'_0}
=\mathcal{J}\mathcal{O}_{Y'}\).
\begin{enumeratei}
\item One has a natural morphism
\[
b_{f_0}\colon
\Hom_{\mathcal{O}_{Y'_0}}\bigl(
i'^*_0 g^*_0(\mathcal{N}_{Y/X}\otimes\mathcal{O}_{Y_0}),
\mathcal{J}\mathcal{O}_{Y'}\bigr)
\longrightarrow
\Hom_{\mathcal{O}_{T_0}}\bigl(
f^*_0 g^*_0(\mathcal{N}_{Y/X}\otimes\mathcal{O}_{Y_0}),
\mathcal{J}\mathcal{O}_T\bigr).
\]
\item One also has a natural morphism, functorial in~\(T\)%
{\renewcommand{\thefootnote}{(100)}\nde{One has deleted the hypothesis
``\(Y'_0\) locally a complete intersection in \(X'_0\)'', superfluous
by~4.8.0; on the other hand, one has added that \(a_{g_0}(f_0)\) is
``functorial in \(T\)'', this playing a crucial role in the proof
of~4.17.}}
\[
a_{g_0}(f_0)\colon
\Hom_{\mathcal{O}_{T_0}}\bigl(
f^*_0(\mathcal{N}_{Y'/X'}\otimes\mathcal{O}_{Y'_0}),
\mathcal{J}\mathcal{O}_T\bigr)
\longrightarrow
\Hom_{\mathcal{O}_{T_0}}\bigl(
f^*_0 g^*_0(\mathcal{N}_{Y/X}\otimes\mathcal{O}_{Y_0}),
\mathcal{J}\mathcal{O}_T\bigr).
\]
\end{enumeratei}
\end{lemma}

\begin{proof}
\renewcommand{\qedsymbol}{}
{\renewcommand{\thefootnote}{(101)}\nde{One has spelled out the proof, in
order to bring out the ``functoriality in \(T\)'' of~\(a_{g_0}\).}}%
Let us first remark that, \(X\), \(X'\), \(Y\), \(Y'\) being fixed, giving
a \(T\) as above is equivalent to giving a morphism
\((f,f_{\mathcal{J}})\colon T\to X'\times_{X'_+}Y'_+\).
Set \(Z=X'\times_{X'_+}Y'_+\) and denote by \(M\) and \(M'\) the
\(Z\)-functors defined by: for every \(f\colon T\to Z\),
\[
\begin{aligned}
M(T)
&=\Hom_{\mathcal{O}_{T_0}}\bigl(
f^*_0 g^*_0(\mathcal{N}_{Y/X}\otimes\mathcal{O}_{Y_0}),
\mathcal{J}\mathcal{O}_T\bigr),\\
M'(T)
&=\Hom_{\mathcal{O}_{T_0}}\bigl(
f^*_0(\mathcal{N}_{Y'/X'}\otimes\mathcal{O}_{Y'_0}),
\mathcal{J}\mathcal{O}_T\bigr).
\end{aligned}
\]
{\renewcommand{\thefootnote}{(102)}\nde{The situation will simplify
from~4.16 onwards: one will restrict oneself to schemes \emph{flat}
over~\(S\), \(Y\) will be a flat \(S\)-group and \(Y'=Y\times_S Y\), and one
will then obtain \(S_0\)-functors \(N_0\) and \(N'_0\).}}%
One has in any case a commutative diagram:
\[
\xymatrix@C=4.5em@R=1.8em{
f_0^*(\mathcal{J}\otimes_{\mathcal{O}_{S_0}}\mathcal{O}_{Y'_0}) \ar[d]
& \mathcal{J}\otimes_{\mathcal{O}_{S_0}}\mathcal{O}_{T_0} \ar[d] \\
f_0^*(\mathcal{J}\mathcal{O}_{Y'}) \ar@{-->}[r]
& \mathcal{J}\mathcal{O}_T
}
\]
and as \(Y'\) is flat over \(S\), the left-hand arrow is an isomorphism,
so one obtains a morphism of \(\mathcal{O}_{T_0}\)-modules
\(f_0^*(\mathcal{J}\mathcal{O}_{Y'})\to\mathcal{J}\mathcal{O}_T\). The
latter induces a morphism of abelian groups
\[
\Hom_{\mathcal{O}_{T_0}}\bigl(
f^*_0 g^*_0(\mathcal{N}_{Y/X}\otimes\mathcal{O}_{Y_0}),
f^*_0(\mathcal{J}\mathcal{O}_{Y'})\bigr)
\longrightarrow M(T)
\]
and, composing with the morphism
\[
M(Y')\longrightarrow
\Hom_{\mathcal{O}_{T_0}}\bigl(
f^*_0 g^*_0(\mathcal{N}_{Y/X}\otimes\mathcal{O}_{Y_0}),
f^*_0(\mathcal{J}\mathcal{O}_{Y'})\bigr),
\]
induced by \(f_0^*\), one obtains the morphism
\(b_{f_0}\colon M(Y')\to M(T)\).

Likewise, one has in any case a diagram
\[
\xymatrix@C=3.2em@R=1.9em{
g_0^*(\mathcal{N}_{Y/X}\otimes_{\mathcal{O}_Y}\mathcal{O}_{Y_0})
  \ar@{-->}[r] \ar[d]
& \mathcal{N}_{Y'/X'}\otimes_{\mathcal{O}_{Y'}}\mathcal{O}_{Y'_0} \ar[d] \\
g_0^*(\mathcal{N}_{Y_0/X_0}) \ar[r]
& \mathcal{N}_{Y'_0/X'_0}
}
\]
and as \(Y'\) is flat over \(S\), the second vertical arrow is an
isomorphism, by~4.8.0. One therefore obtains an
\(\mathcal{O}_{Y'_0}\)-morphism
\[
i'^*_0 g^*_0(\mathcal{N}_{Y/X}\otimes_{\mathcal{O}_Y}\mathcal{O}_{Y_0})
\longrightarrow
\mathcal{N}_{Y'/X'}\otimes_{\mathcal{O}_{Y'}}\mathcal{O}_{Y'_0}
\]
which induces a morphism
\(a_{g_0}(\mathrm{id}_{Y'})\colon M'(Y')\to M(Y')\) and, for every
\(f\colon T\to Z\), a morphism
\(a_{g_0}(f)\colon M'(T)\to M(T)\) such that one has a commutative diagram
\[
\xymatrix@C=4.2em@R=1.7em{
M'(Y') \ar[r]^{a_{g_0}(\mathrm{id}_{Y'})} \ar[d]_{b'_{f_0}}
& M(Y') \ar[d]^{b_{f_0}} \\
M'(T) \ar[r]_{a_{g_0}(f_0)}
& M(T)
}
\]
(where \(b'_{f_0}\) is defined as \(b_{f_0}\)).
\hfill Q.E.D.
\end{proof}

\begin{remark}{4.12.1}\label{III.4.12.1}
{\renewcommand{\thefootnote}{(103)}\nde{One has added this remark, used
in the proof of~4.17.}}%
Let us denote by \(M_0\) and \(M'_0\) the \(Y'_0\)-functors defined by: for
every \(f\colon T_0\to Y'_0\),
% typo?: printed M_0(T) and M'_0(T), while f\colon T_0\to Y'_0; kept as printed
\[
\begin{aligned}
M_0(T)
&=\Hom_{\mathcal{O}_{T_0}}\bigl(
f^*_0 g^*_0(\mathcal{N}_{Y/X}\otimes\mathcal{O}_{Y_0}),
\mathcal{J}\otimes_{\mathcal{O}_{S_0}}\mathcal{O}_{T_0}\bigr),\\
M'_0(T)
&=\Hom_{\mathcal{O}_{T_0}}\bigl(
f^*_0(\mathcal{N}_{Y'/X'}\otimes\mathcal{O}_{Y'_0}),
\mathcal{J}\otimes_{\mathcal{O}_{S_0}}\mathcal{O}_{T_0}\bigr).
\end{aligned}
\]
Let us remark at once that \(Z_0=Y'_0\) and that on the category of
\(Z\)-schemes \(T\) which are \emph{flat} over~\(S\), \(M\) and \(M'\)
coincide, respectively, with the functors
\(\prod_{S_0/S}M_0\) and \(\prod_{S_0/S}M'_0\). In this case, \(b_{f_0}\)
is simply the morphism
% typo?: codomain printed M(T_0), not M_0(T_0); kept as printed
\[
f_0^*\colon M_0(Y'_0)\longrightarrow M(T_0)
\]
induced by~\(f_0\), and for every morphism \(u\colon U\to T\), setting
\(h=f\circ u\), one has a commutative diagram
\[
\xymatrix@C=4em@R=1.7em{
M'_0(T_0) \ar[r]^{a_{g_0}(f_0)} \ar[d]_{u_0^*}
& M_0(T_0) \ar[d]^{u_0^*} \\
M'_0(U_0) \ar[r]_{a_{g_0}(h_0)}
& M_0(U_0)
}
\]
\ie \(a_{g_0}\) becomes a morphism of functors
\(\prod_{S_0/S}M'_0\to\prod_{S_0/S}M_0\).
\end{remark}

\begin{proposition}{4.13}\label{III.4.13}
Suppose \(Y'\) flat over \(S\). One then has the formula:
\[
c(X,Y,g\circ f)
=a_{g_0}\bigl(c(X',Y',f)\bigr)
+b_{f_0}\bigl(c(X,Y,g\circ i')\bigr).
\]
\end{proposition}

As the definition of the various obstructions and of the morphisms
\(a_{g_0}\) and \(b_{f_0}\) is local, one sees easily that it suffices to
check the given formula when the various schemes in question are affine.
Let us therefore write \(S=\Spec(\Lambda)\), \(S_J=\Spec(\Lambda/J)\),
\(S_0=\Spec(\Lambda/I)\), \(T=\Spec(C)\), \(X=\Spec(A)\),
\(Y=\Spec(A/I_Y)=\Spec(B)\), \(X'=\Spec(A')\),
\(Y'=\Spec(A'/I_{Y'})=\Spec(B')\).

One therefore has a diagram of rings and ideals%
{\renewcommand{\thefootnote}{(104)}\nde{One has kept the notations of the
original, writing \(f\colon A'\to C\) and \(g\colon A\to A'\) for the ring
morphisms corresponding to \(f\colon T\to X'\) and \(g\colon X'\to X\). This
explains the formula \(c(X,Y,g\circ f)(x\otimes 1)=f(g(x))\), for
\(x\in I_Y\).}}
\[
\xymatrix@C=3em@R=1.55em{
& B' & B \\
C & A' \ar[l]_{f} \ar[u]_{\pi'} & A \ar[l]_{g} \ar[u]_{\pi} \\
& I_{Y'} \ar@{^{(}->}[u] & I_Y \ar@{^{(}->}[u]\,.
}
\]
Let us study the various terms of the formula to be proved. In what
follows, if \(x\in I_Y\) (\resp \(u\in I_{Y'}\)), one denotes by \(x\)
(\resp \(u\)) its image in \(I_Y/I_Y^2\)
(\resp \(I_{Y'}/I_{Y'}^2\)); on the other hand, if \(m\) belongs to a
\(\Lambda\)-module \(M\), one denotes by \(m_0\) its image in
\(M_0=M/IM\).

One has seen that \(c=c(X,Y,g\circ f)\) is the unique \(C_0\)-morphism
\(I_Y/I_Y^2\otimes_B C_0\to JC\) such that \(c(x\otimes 1)=f(g(x))\), for
every \(x\in I_Y\).

Let us fix \(x\in I_Y\); one has \(g(x)\in I_{Y'}+JA'\) since
\(g_J(Y'_J)\subset Y_J\). Let us write
\(g(x)=x'+\sum\lambda_i a'_i\), with \(x'\in I_{Y'}\), \(\lambda_i\in J\),
\(a'_i\in A'\). One therefore has
\begin{equation}
c(X,Y,g\circ f)(x\otimes 1)
=f(g(x))=f(x')+\sum\lambda_i f(a'_i).
\tag{1}
\end{equation}

Let us now consider \(a_{g_0}(c(X',Y',f))\). By the definitions set down,
it is defined by the diagram
\[
\xymatrix@C=1.5em@R=2.15em{
& I_{Y'} \ar[d] \ar[r]^{f} & C \\
I_{Y'_0}/I_{Y'_0}^2\otimes_{B'_0}C_0
& I_{Y'}/I_{Y'}^2\otimes_{B'}C_0
  \ar[l]^{\sim}
  \ar[r]^{c(X',Y',f)}
& JC \\
I_{Y_0}/I_{Y_0}^2\otimes_{B_0}C_0 \ar[u]_{g_0}
& I_Y/I_Y^2\otimes_B C_0
  \ar[l]
  \ar[u]
  \ar[ur]_{a_{g_0}(c(X',Y',f))}
&
}
\]
One therefore has
\(a_{g_0}(c(X',Y',f))(x\otimes 1)=f(u)\), where \(u\) is an element of
\(I_{Y'}\) whose image \(u\) in \(I_{Y'_0}/I_{Y'_0}^2\) satisfies
% typo?: the two factors g_0(x_0) are printed alike; kept as printed
\(u_0\otimes 1=g_0(x_0)\otimes 1=g_0(x_0)\otimes 1\).
One may therefore take \(u=x'\), and one finds
\begin{equation}
a_{g_0}(c(X',Y',f))(x\otimes 1)=f(x').
\tag{2}
\end{equation}

Let us finally consider \(b_{f_0}(c(X,Y,g\circ i'))\). By hypothesis, the
morphism of \(\Lambda_0\)-algebras \(f_0\colon A'_0\to C_0\) factors through
\(B'_0\) and hence, as
\(J\otimes_{\Lambda_0}B'_0\xrightarrow{\sim}JB'\)
(\(B'\) being flat over~\(\Lambda\)), one obtains a morphism of
\(B'_0\)-modules \(\psi\colon JB'\to JC\) such that one has a commutative
diagram:
\[
\xymatrix@C=2.6em@R=1.8em{
J\otimes_{\Lambda_0}A'_0 \ar[r]^{\mathrm{id}\otimes\pi'} \ar[d]
& J\otimes_{\Lambda_0}B'_0 \ar[r]^{\mathrm{id}\otimes f_0} \ar[d]^{\wr}
& J\otimes_{\Lambda_0}C_0 \ar[d] \\
JA' \ar[r]_{\pi'}
& JB' \ar@{-->}[r]_{\psi}
& JC.
}
\]
Denote by \(\phi\colon JB'\otimes_{B'_0}C_0\to JC\) the morphism of
\(C_0\)-modules deduced from~\(\psi\); then one has, for every
\(a'\in A'\), \(\lambda\in J\),
\begin{equation}
\phi(\lambda\pi'(a')\otimes 1)=\lambda f(a').
\tag{\dag}
\end{equation}

Then, \(b_{f_0}(c(X,Y,g\circ i'))\) is defined by the commutative diagram:
\[
\xymatrix@C=2.4em@R=1.85em{
I_Y \ar[rr]^{\pi'\circ g} \ar[d]
& & B' \\
I_Y/I_Y^2\otimes_B B'_0
  \ar[rr]^{c(X,Y,g\circ i')} \ar[d]
& & JB' \ar[d] \\
I_Y/I_Y^2\otimes_B C_0
  \ar[rr] \ar[drr]_{b_{f_0}(c(X,Y,g\circ i'))}
& & JB'\otimes_{B'_0}C_0 \ar[d]^{\phi} \\
& & JC.
}
\]
One therefore has at once
\begin{equation}
b_{f_0}(c(X,Y,g\circ i'))(x\otimes 1)
=\phi\Bigl(\sum\lambda_i\pi'(a'_i)\otimes 1\Bigr)
=\sum\lambda_i f(a'_i),
\tag{3}
\end{equation}
the last equality following from~\((\dagger)\) above. The comparison of
the three explicit results~(1), (2), (3) gives the desired formula.

\begin{corollary}{4.14}\label{III.4.14}
Let \(Y\), \(Y'\) be two subschemes of \(X\), flat, reducing
to~\(Y_{\mathcal{J}}\); suppose \(Y_0\) locally a complete intersection
in~\(X_0\). If \(f\colon T\to X\) is an \(S\)-morphism such that
\(f_{\mathcal{J}}\) factors through
\(Y_{\mathcal{J}}\to X_{\mathcal{J}}\), one has the formula
\[
c(X,Y,f)-c(X,Y',f)=b_{f_0}\bigl(d(Y,Y')\bigr).
\]
\end{corollary}

Indeed, applying the preceding formula to the diagram
\[
\xymatrix@C=3.2em@R=1.6em{
& Y' \ar@{^{(}->}[d]^{i'} & Y \ar@{^{(}->}[d]^{i} \\
T \ar[r]^{f} & X \ar[r]^{\mathrm{id}} & X
}
\]
one finds \(c(X,Y,f)-c(X,Y',f)=b_{f_0}(c(X,Y,i'))\). Moreover, by~4.8\,(ii),
one has \(c(X,Y,i')=d(Y,Y')\).

\begin{proposition}{4.15}\label{III.4.15}
Let \(X\) be an \(S\)-group smooth over \(S\) and \(Y\) an \(S\)-subgroup flat
and locally of finite presentation over \(S\). Then \(Y\) is locally a
complete intersection (\cf 4.6.2) in~\(X\).
\end{proposition}

\begin{proof}
{\renewcommand{\thefootnote}{(105)}\nde{One has added to the statement the
hypothesis that \(Y\) be locally of finite presentation over \(S\), and one
has given the proof which follows, more direct than the one sketched in
the original. For completeness, let us also spell out the latter. As in
the proof given above, one first reduces to the case where
\(S=\Spec(k)\), \(k\) being an algebraically closed field. By
EGA~IV\(_4\), 16.9.10 and~19.3.2, it suffices to see that, for every
\(y\in Y\), the completion of the local ring \(\mathcal{O}_{Y,y}\) is the
quotient of a complete noetherian local ring by a regular sequence. By
\loccit, 19.3.3, the set of \(y\in Y\) satisfying this property is an open
subset \(U\) of \(Y\); as \(Y\) is of finite type over \(k\), it suffices
to show that \(U\) contains every closed point. As \(Y\) is a \(k\)-group
it suffices, by a translation argument, to show that the property holds
for the completion of \(\mathcal{O}_{Y,e}\), \ie for the ``formal group''
\(\widehat{Y}\) corresponding to \(Y\) (\cf Exp.~\(\mathrm{VII}_{\mathrm{B}}\)).
Now, as \(X\) is smooth, the affine algebra \(A(\widehat{X})\) is an algebra
of formal power series \(k[[X_1,\ldots,X_n]]\), and one concludes by means
of the structure theorem of Dieudonn\'e, which shows that \(A(\widehat{Y})\)
is isomorphic to a quotient
\(k[[X_1,\ldots,X_{r+s}]]/(X_1^{p^{n_1}},\ldots,X_r^{p^{n_r}})\),
where \(p\) is the characteristic exponent of \(k\) and
\(r+s\leqslant n\), \cf \(\mathrm{VII}_{\mathrm{B}}\), Remark~5.5.2\,(b).}}%
One will show that the immersion \(Y\to X\) is regular in the sense of
EGA~IV\(_4\), 16.9.2, which implies that it is so also in the sense
of~4.6.2, by EGA~IV\(_4\), 19.5.1 (moreover, by \loccit, the two
definitions are equivalent if \(S\) is locally noetherian). Thus, in what
follows, one takes ``regular immersion'' in the sense of
EGA~IV\(_4\), 16.9.2. As \(X\) and \(Y\) are flat and locally of finite
presentation over \(S\), then, by EGA~IV\(_4\), 19.2.4, it suffices to
show that, for every \(s\in S\), \(Y_s\to X_s\) is a regular immersion. By
EGA~IV\(_4\), 19.1.5\,(ii), one is reduced to checking the assertion on
the geometric fibers of \(S\), hence when \(S\) is the spectrum of an
algebraically closed field \(k\).

Then, by \(\mathrm{VI}_{\mathrm{A}}\), 3.2, the quotient \(X/Y\) exists
and is smooth, the morphism \(\pi\colon X\to X/Y\) is flat, and one has a
cartesian square
\[
\xymatrix@C=3.2em@R=1.6em{
Y \ar[r]^{f} \ar[d] & X \ar[d]^{\pi} \\
e \ar[r]_{i} & X/Y
}
\]
(where \(e\) is the image in \(X/Y\) of the unit point of \(X\)). Thus, by
flat base change (\cf EGA~IV\(_4\), 19.1.5\,(ii)), it suffices to see that
\(i\) is a regular immersion, which is immediate since the noetherian local
ring \(\mathcal{O}_{X/Y,e}\) is smooth, hence its maximal ideal is
generated by a regular sequence.
\end{proof}

\numpara{4.16}\label{III.4.16}
{\renewcommand{\thefootnote}{(106)}\nde{One has reorganized~4.16 by
gathering in it, on the one hand, the hypotheses stated at the end of~4.15
and, on the other hand, the definition of the obstruction \(DY\).}}%
Let \(X\) be an \(S\)-group \emph{smooth} over \(S\); denote by
\(\mu\colon X\times_S X\to X\) its group law.
Let us be given an \(S_{\mathcal{J}}\)-subgroup \(Y_{\mathcal{J}}\) of
\(X_{\mathcal{J}}\), \emph{flat and locally of finite presentation}
over~\(S_{\mathcal{J}}\). By~4.15, \(Y_{\mathcal{J}}\) is locally a
complete intersection in~\(X\).
% typo?: kept as printed (``dans X'', not ``dans X_J'')

Thus, by~4.6.5, every \emph{flat}%
{\renewcommand{\thefootnote}{(107)}\nde{One has corrected the original by
adding ``flat''.}}
\(S\)-scheme \(Y\) lifting \(Y_{\mathcal{J}}\) is locally a complete
intersection in~\(X\). For such a \(Y\) one has, by~4.8.0,
\begin{equation}
\mathcal{N}_{Y/X}\otimes_{\mathcal{O}_Y}\mathcal{O}_{Y_0}
=\mathcal{N}_{Y_0/X_0}
=\mathcal{N}_{Y_{\mathcal{J}}/X_{\mathcal{J}}}
  \otimes_{\mathcal{O}_{Y_{\mathcal{J}}}}\mathcal{O}_{Y_0}.
\tag{4.16.1}
\end{equation}

On the other hand, let us denote by \(\varepsilon_0\colon S_0\to Y_0\) the
unit section of \(Y_0\) and by \(\mathfrak{n}_{Y_0/X_0}\) the
quasi-coherent \(\mathcal{O}_{S_0}\)-module:
\[
\mathfrak{n}_{Y_0/X_0}=\varepsilon_0^*(\mathcal{N}_{Y_0/X_0}).
\]
As \(Y_0\) and \(X_0\) are \(S_0\)-groups, one sees easily that
\(\mathcal{N}_{Y_0/X_0}\) is invariant under the translations (say on the
left) of \(Y_0\), hence%
{\renewcommand{\thefootnote}{(108)}\nde{See~4.25 below.}}
is the inverse image under \(Y_0\to S_0\) of \(\mathfrak{n}_{Y_0/X_0}\),
\ie one has
\begin{equation}
\mathcal{N}_{Y_0/X_0}
=\mathfrak{n}_{Y_0/X_0}\otimes_{\mathcal{O}_{S_0}}\mathcal{O}_{Y_0}.
\tag{4.16.2}
\end{equation}
Taking~(4.16.1) and~(4.16.2) into account, one deduces on the one hand
from~4.5 that the set of sub-\(S\)-schemes \(Y\) of \(X\), \emph{flat}
over~\(S\), lifting \(Y_{\mathcal{J}}\), is empty or a principal
homogeneous space under
\begin{equation}
\Hom_{\mathcal{O}_{Y_0}}\bigl(
\mathfrak{n}_{Y_0/X_0}\otimes_{\mathcal{O}_{S_0}}\mathcal{O}_{Y_0},\,
\mathcal{J}\otimes_{\mathcal{O}_{S_0}}\mathcal{O}_{Y_0}\bigr),
\tag{4.16.3}
\end{equation}
and one deduces on the other hand from~4.8\,(i) that, for every such \(Y\)
and every \(S\)-morphism \(f\colon T\to X\) such that
\(f_{\mathcal{J}}\colon T_{\mathcal{J}}\to X_{\mathcal{J}}\) factors
through \(Y_{\mathcal{J}}\), the obstruction \(c(X,Y,f)\) to \(f\)
factoring through \(Y\) is an element of
\[
\Hom_{\mathcal{O}_{T_0}}\bigl(
\mathfrak{n}_{Y_0/X_0}\otimes_{\mathcal{O}_{S_0}}\mathcal{O}_{T_0},\,
\mathcal{J}\mathcal{O}_T\bigr);
\]
if moreover \(T\) is \emph{flat} over \(S\), this last group equals
\[
\Hom_{\mathcal{O}_{T_0}}\bigl(
\mathfrak{n}_{Y_0/X_0}\otimes_{\mathcal{O}_{S_0}}\mathcal{O}_{T_0},\,
\mathcal{J}\otimes_{\mathcal{O}_{S_0}}\mathcal{O}_{T_0}\bigr).
\]
This leads one to introduce the group functor \(N_0\) below:

\begin{definition}{4.16.1}\label{III.4.16.1}
Let \(N_0\) be the commutative group \(S_0\)-functor defined by: for every
\(Z\in\Ob\Sch/S_0\),
\begin{equation}
\Hom_{S_0}(Z,N_0)
=\Hom_{\mathcal{O}_Z}\bigl(
\mathfrak{n}_{Y_0/X_0}\otimes_{\mathcal{O}_{S_0}}\mathcal{O}_Z,\,
\mathcal{J}\otimes_{\mathcal{O}_{S_0}}\mathcal{O}_Z\bigr).
\tag{*}
\end{equation}
Then the set of sub-\(S\)-schemes \(Y\) of \(X\), \emph{flat} over \(S\),
lifting \(Y_{\mathcal{J}}\), is empty or a principal homogeneous space under
\[
\Hom_{S_0}(Y_0,N_0)=C^1(Y_0,N_0).
\]
For each such \(Y\), let us consider the following diagram:
\[
\xymatrix@C=3.4em@R=1.6em{
Y\times_S Y \ar@{^{(}->}[d]_{(i,i)} & Y \ar@{^{(}->}[d]^{i} \\
X\times_S X \ar[r]_{\mu} & X
}
\]
and let us denote by \(DY=c(X,Y,\mu\circ(i,i))\) the obstruction to
\(\mu\circ(i,i)\) factoring through \(Y\), \ie to \(Y\) being stable under
the group law of \(X\); by what precedes, \(DY\) is an element of
\[
N_0(Y_0\times_{S_0}Y_0)=C^2(Y_0,N_0).
\]
\end{definition}
